\begin{helper}
Let $C$ be a nonempty finite set of candidates, let $p\in[0,1]$, and assign
each candidate $x\in C$ a priority $\pi(x)\in[0,1]$.  Then there is a
candidate $x_0\in C$ with maximal priority on $C$, and
\[
1-p(1-\pi(x_0))
\le
1-\prod_{x\in C} p(1-\pi(x)).
\]
\end{helper}
\begin{proof}
Since $C$ is finite and nonempty, the finite set of real priorities
$\{\pi(x):x\in C\}$ has a maximum.  Choose $x_0\in C$ attaining this maximum,
so $\pi(x)\le \pi(x_0)$ for every $x\in C$.

For every $x\in C$, the assumptions $0\le p\le1$ and
$0\le \pi(x)\le1$ imply $0\le p(1-\pi(x))\le1$.  Applying
\noderef{SamplingOneBlockerProductInequality} to the finite set $C$, the
activation probability $p$, the priority function $\pi$, and the chosen
element $x_0$ gives
\[
\prod_{x\in C}p(1-\pi(x))\le p(1-\pi(x_0)).
\]
Subtracting the two sides from $1$ reverses the displayed comparison exactly
as claimed:
\[
1-p(1-\pi(x_0))
\le
1-\prod_{x\in C}p(1-\pi(x)).
\]
\end{proof}
